\documentclass[11pt,letterpaper]{article}
\usepackage[T1]{fontenc}
\usepackage{lmodern}
\usepackage[margin=0.93in]{geometry}
\usepackage{amsmath,amssymb,amsthm,mathtools}
\usepackage{microtype,booktabs,longtable,tabularx,array}
\usepackage{xcolor,graphicx,enumitem,fancyhdr}
\usepackage{needspace}
\usepackage{listings}
\usepackage[unicode,colorlinks=true,linkcolor=blue!45!black,citecolor=teal!55!black,
 urlcolor=blue!55!black,breaklinks=true]{hyperref}
\definecolor{ink}{HTML}{16324F}
\definecolor{accent}{HTML}{126782}
\definecolor{light}{HTML}{F2F6FA}
\hypersetup{pdftitle={Toward Quasi-Polynomial Unknot Recognition},
 pdfauthor={Research continuation for the ProveIt project}}
\pagestyle{fancy}
\fancyhf{}
\fancyhead[L]{\small\color{ink}ProveIt: Unknot Recognition}
\fancyhead[R]{\small October 2026}
\fancyfoot[C]{\thepage}
\setlength{\headheight}{14pt}
\setlength{\emergencystretch}{2em}
\makeatletter
\renewcommand*\l@subsection{\@dottedtocline{2}{1.5em}{3.1em}}
\makeatother
\setlist{itemsep=3pt,topsep=5pt}
\lstset{basicstyle=\ttfamily\small,breaklines=true,columns=fullflexible,
 frame=single,rulecolor=\color{black!20},backgroundcolor=\color{light},
 showstringspaces=false,keepspaces=true}
\newtheorem{theorem}{Theorem}[section]
\newtheorem{lemma}[theorem]{Lemma}
\newtheorem{proposition}[theorem]{Proposition}
\newtheorem{corollary}[theorem]{Corollary}
\theoremstyle{definition}
\newtheorem{definition}[theorem]{Definition}
\newtheorem{remark}[theorem]{Remark}
\newcommand{\F}{\mathbb F_2}
\newcommand{\Z}{\mathbb Z}
\newcommand{\N}{\mathbb N}
\newcommand{\SL}{\operatorname{SL}}
\newcommand{\PSL}{\operatorname{PSL}}
\newcommand{\tr}{\operatorname{tr}}
\newcommand{\Kh}{\operatorname{Kh}}
\newcommand{\Cat}{\operatorname{Cat}}
\newcommand{\Hom}{\operatorname{Hom}}
\newcommand{\End}{\operatorname{End}}
\newcommand{\rank}{\operatorname{rank}}
\newcommand{\softO}{\widetilde O}
\newcommand{\code}[1]{\texttt{\detokenize{#1}}}
\title{\color{ink}\bfseries Toward Quasi-Polynomial\\Unknot Recognition\\[5pt]
 \large An exact three-braid backend with linear symbolic cost,\\
 certified Markov descent, and intrinsic scanner multiplicities}
\author{Research continuation for Vladimir Reshetnikov's ProveIt project}
\date{7 October 2026}
\begin{document}
\maketitle
\begin{abstract}
We extend ProveIt's exact unknot-recognition implementation,
\code{fastunknot}.
A supplied three-strand braid can be decided by reducing a word in
\(C_2*C_3\), retaining its exponent sum to resolve the central ambiguity.
The resulting complete specialized algorithm uses \(O(\ell)\) symbol
operations on an explicit word of length \(\ell\); a conservative bit bound
is \(O(\ell\log(\ell+2))\). An independent integer-matrix implementation
checks the same classification in \(O(\ell^2)\) elementary bit operations.
We integrate these procedures with a general braid-writhe obstruction,
replayable elementary Markov reductions, and the existing exact fallback.

We also identify the object multiplicities that remain after exhaustive
categorical cancellation: they are Betti numbers of a scalar residue
complex, and are independent of pivot choices at a fixed prefix.
A diagnostic verifies this identity against the actual scanning engines.
For disk frontiers, a nilpotent radical gives a finite homological
perturbation formula for batch cancellation. We derive explicit
width--multiplicity complexity bounds and explain why they are conditional
on a separate multiplicity bound.

The topological classification and algebraic tools used here are established
mathematics. The contributions are their convention-checked specialization,
implementation, certification, cost accounting, and experimental evaluation.
We do not obtain a general \(n^{O(\log n)}\) recognizer. The report locates
the remaining obligations in the Khovanov and geometric-hierarchy routes,
taking account of both relevant 2026 preprints.
\end{abstract}
\tableofcontents
\newpage

\section{Outcome, scope, and the exact baseline}
\label{sec:scope}

The concrete improvement is a complete inexpensive decision route for a
verified low-strand braid representation, before the general pipeline has
to build a Khovanov complex. On other braid inputs we add two certified
opportunities to finish early: a writhe obstruction and explicit endpoint
destabilizations. For planar-diagram and grid inputs, the existing exact
recognizer remains available. We do not infer a low braid index from a
small scan frontier.

The baseline is the \code{fast/} subtree at ProveIt commit
\begin{center}
\small\code{4e6fe879e5238ec0b134b6d33fdca0b8c7c1c711}.
\end{center}
Its Git tree is
\code{b99a666269f660d932fc82ac8d1f29d3ae42231f}.
The last commit affecting the larger unknot-recognition subtree, when
inspected, was \code{78333302d9babf0efb1a5d218e122f13c35bb306}.
The supplied provenance file records the source and hashes; this is a
continuation of the integrated implementation, not a comparison against
the earlier 0.1 prototype.

That distinction matters. The current baseline already contains
Reidemeister I/II reduction, a bounded Reidemeister III search, descending
tests, modular Alexander and bracket filters, visible connected-sum
factorization, improved scan ordering, bit-packed cobordisms, interned
matchings, cached gluing plans, and several pivot and racing strategies.
Repeating those optimizations would not address the next structural
bottleneck. The preceding nine proposals and synthesis are acknowledged
as prior project work \cite{proveit}.

\begin{center}
\begin{tabularx}{\textwidth}{@{}>{\raggedright\arraybackslash}p{0.29\textwidth}
 >{\raggedright\arraybackslash}X >{\raggedright\arraybackslash}p{0.22\textwidth}@{}}
\toprule
Result & Precise scope & Status in this bundle\\
\midrule
Three-braid decision & Explicit validated braid word on at most three strands
 & Proved and implemented\\
General writhe rejection & Unknotted \(m\)-braids satisfy \(|e|\le m-1\)
 & Proved from a cited theorem; implemented\\
Endpoint Markov descent & A sufficient, replayable destabilization criterion
 & Proved and implemented\\
Canonical multiplicities & A fixed prefix and a fixed homological normalization
 & Proved and checked on actual complexes\\
Finite radical batching & Disk arc algebra over \(\F\)
 & Proved; batching implementation pending\\
General quasi-polynomial decision & Arbitrary classical knot diagrams
 & Not achieved by this work\\
\bottomrule
\end{tabularx}
\end{center}

\subsection{What is and is not a novelty claim}

Polynomial algorithms for three-braid conjugacy, the classification of
closed three-braids, and polynomial special-purpose computations are
classical or already in the literature
\cite{birmanmenasco,datta2020,datta2022,kelomaki2026}.
The residue and perturbation arguments specialize established homological
algebra \cite{khovanovtangles,crainic}.
We present full proofs of the specific criteria and bounds needed by this
implementation, rather than claim discovery of those general theories.

The new project-level outcome is an executable, independently checked route
that avoids explicit homological bases on a complete input class, together
with a diagnostic that tells us which multiplicities pivot tuning cannot
remove. These are concrete performance improvements and a better specified
research program. Their value does not depend on attaching an unsupported
claim of a general breakthrough.

\subsection{Input and computational model}

A braid word is an explicit list
\(\beta=\sigma_{|g_1|}^{\operatorname{sgn}(g_1)}
\cdots\sigma_{|g_\ell|}^{\operatorname{sgn}(g_\ell)}\)
with \(1\le |g_j|<m\).
The permutation obtained by exchanging adjacent strands is checked to be
one cycle. Thus the input describes a knot, not a general link.
For a one-component closure, \(\ell\ge m-1\): every adjacent gap in the
strand path must be crossed at least once for the permutation to be
transitive. An empty word on one strand is the crossing-free unknot.

A symbol operation acts on a generator, a free-product token, or a
logarithmic-size index. A conservative bit interpretation charges
\(O(\log(\ell+m+2))\) for such an operation. We state separately the
larger-integer costs of the matrix method. The bounds do not apply to a
binary exponent or straight-line program without charging for its
expansion. No succinct-input assumption is hidden in \(\ell\).

The direct word API has the strongest stated specialized bound. Conversion
to a PD code and its existing validation are separate preprocessing in the
public \code{Diagram} API. Keeping braid provenance allows the public
recognizer to use the word route after parsing; it does not make all
existing parser operations constant time.

\section{The relevant state of the art}
\label{sec:literature}

There are two different reasons to improve this implementation: obtain
better exact decisions now, and develop a route to a stronger general
complexity theorem. They should not be conflated.

Lackenby's 2021 announcement describes a quasi-polynomial unknot
algorithm \cite{lackenbytalk}. The July 2026 preprint supplies a hierarchy
algorithm and additional constructive and certificate results, but its
Section 9 separates iteration counts from the cost of its individual
steps \cite{lackenby2026}. Our project has not implemented the announced
general bound. We also do not infer from the absence of a cost proof in
that preprint that the announcement is false.

The January 2026 preprint of Kelom\"aki and Sch\"utz proves that the
usual scanning and divide-and-conquer Khovanov algorithms can take
exponential time even on positive or alternating three-braids. Their
specialized three-braid computation avoids that obstruction using
structural formulas \cite{kelomaki2026}. This supports a change of
representation or decision problem, not a further claim that width alone
controls the existing explicit complex.

Our implementation asks only whether the knot is trivial. It need not
compute every homology group, and it need not construct a minimal chain
complex whose number of basis elements can itself be exponential.
The three-braid route is a particularly clean example of this separation:
an exact decision reduces to a short group-theoretic conjugacy test.

\section{A complete three-braid classifier by free-product reduction}
\label{sec:three}

Write \(a=\sigma_1\), \(b=\sigma_2\), and use the standard presentation
\[
 B_3=\langle a,b\mid aba=bab\rangle.
\]
Put
\[
 x=aba,\qquad y=ab,\qquad z=x^2=y^3=(ab)^3.
\]
Then
\[
 B_3=\langle x,y\mid x^2=y^3\rangle,\qquad
 Z(B_3)=\langle z\rangle,\qquad
 B_3/\langle z\rangle\cong C_2*C_3.
\]
These familiar presentations can also be checked by substituting
\(a=y^{-1}x\), \(b=x^{-1}y^2\) in both directions.
Let \(e:B_3\to\Z\) be the exponent-sum homomorphism; \(e(z)=6\).

\subsection{The topological input}

We use the following consequence of the closed-three-braid classification
\cite{birmanmenasco}; Datta states the unknot fiber explicitly
\cite[Theorem 4.1(i)]{datta2022}.

\begin{theorem}[Classical three-braid unknot classification]
\label{thm:classical}
The braids in \(B_3\) with unknot closure form precisely the conjugacy
classes represented by
\[
 ab,\qquad (ab)^{-1},\qquad ab^{-1}.
\]
\end{theorem}

The representatives close to unknots: their two crossings are successive
positive or negative stabilizations of the one-strand unknot, with mixed
signs in the last case. The nontrivial content of the cited theorem is
necessity. We use it as an explicit topological input, rather than
silently assuming that every unknotted braid destabilizes monotonically.

\begin{lemma}[The exponent sum fixes the central lift]
\label{lem:lift}
If two braids have conjugate images in \(C_2*C_3\) and equal exponent
sums, then they are conjugate in \(B_3\).
\end{lemma}
\begin{proof}
Lift a quotient conjugator to \(g\in B_3\). If the two braids are
\(\beta,\gamma\), then
\(g\beta g^{-1}=\gamma z^k\) for some \(k\in\Z\).
Taking exponent sums gives \(e(\beta)=e(\gamma)+6k\).
Equality forces \(k=0\).
\end{proof}

\subsection{An explicit constant-alphabet algorithm}

In the quotient, the four input letters become
\begin{equation}
\label{eq:images}
 a\mapsto y^2x,\quad a^{-1}\mapsto xy,\quad
 b\mapsto xy^2,\quad b^{-1}\mapsto yx.
\end{equation}
These are quotient identities. Replacing \(y^{-1}\) by \(y^2\) in
\(B_3\) itself would lose a central power; that is why the exponent sum
must be retained.

Represent \(x,y,y^2\) by three tokens. Scan at most \(2\ell\) tokens into
a stack. Adjacent \(x\)'s cancel. Adjacent powers of \(y\) combine by
addition modulo three, with exponent zero deleted. If the incoming token
belongs to a different free factor from the top token, push it.
The resulting word alternates between the factors.

For cyclic reduction, use a deque. If the first and last syllables belong
to the same factor, write the word \(u\,v\,w\), with those endpoint
syllables \(u,w\). Conjugation by \(u^{-1}\) changes it into
\(v(wu)\). Delete the endpoints and append their nonidentity product.
This decreases the number of syllables. Repeat until the length is at
most one or the endpoint factors differ. No search over conjugators is
necessary.

\begin{lemma}[The required free-product conjugacy tests]
\label{lem:conj}
In \(C_2*C_3\), an element is conjugate to \(y\) or \(y^2\) exactly when
cyclic reduction yields that one syllable. It is conjugate to
\(y^2xyx\) exactly when cyclic reduction has four alternating syllables,
with one \(y\) and one \(y^2\).
\end{lemma}
\begin{proof}
A reduced free-product word is unique. The usual conjugacy argument can
be seen directly by cancelling a reduced conjugator against the two ends
of a cyclically reduced word. For length at least two, a conjugate that is
again cyclically reduced can only cyclically permute its syllables.
For a word of length one, conjugacy stays within its factor; the factors
here are abelian. The possible four-syllable words in the statement are
exactly the cyclic rotations of \(y^2xyx\). This proves both claims.
\end{proof}

\Needspace{11\baselineskip}
\begin{theorem}[Complete symbolic criterion]
\label{thm:symbolic}
For a validated three-braid word \(\beta\), let \(w\) be the cyclically
reduced quotient word obtained above. Its closure is the unknot if and
only if one of the following holds:
\begin{center}
\begin{tabular}{@{}cll@{}}
\toprule
\(e(\beta)\) & \(w\) & Representative in \(B_3\)\\
\midrule
\(2\) & \(y\) & \(ab\)\\
\(-2\) & \(y^2\) & \((ab)^{-1}\)\\
\(0\) & length four, with \(y\) and \(y^2\) once each & \(ab^{-1}\)\\
\bottomrule
\end{tabular}
\end{center}
\end{theorem}
\begin{proof}
The images of the three representatives are respectively \(y\), \(y^2\),
and \(y^2xyx\), and their exponent sums are \(2,-2,0\).
Necessity follows from Theorem~\ref{thm:classical}.
For sufficiency, Lemma~\ref{lem:conj} gives a quotient conjugacy to the
corresponding representative, and Lemma~\ref{lem:lift} lifts it to a
braid conjugacy. Taking closures preserves its knot type.
\end{proof}

\begin{theorem}[Cost of the direct word backend]
\label{thm:linear}
The symbolic criterion, including word and component validation for at
most three strands, uses \(O(\ell)\) symbol operations and
\(O(\ell)\) stored tokens. A conservative logarithmic-address bit bound is
\(O(\ell\log(\ell+2))\) time and \(O(\ell\log(\ell+2))\) bits of storage.
\end{theorem}
\begin{proof}
Each input generator produces exactly two quotient tokens. A token is
pushed at most once, and every stack cancellation removes a previously
stored token. Thus the total number of stack updates is linear.
Each cyclic-end operation decreases the deque length, so their total is
also linear. The exponent sum has \(O(\log(\ell+2))\) bits.
On at most three strands, component validation maintains a permutation of
constant size. Checking the accepted word types is constant time after
reduction, aside from emitting an optional word witness of length
\(O(\ell)\). Charging logarithmic index and counter costs gives the stated
conservative bit bounds.
\end{proof}

This is a bound for \emph{decision from a supplied braid word}. It neither
constructs a three-braid presentation of an arbitrary PD diagram nor
computes the entire Khovanov homology of the closure.

\section{An independent integer-matrix criterion}
\label{sec:matrix}

An independently structured implementation is valuable here: a sign error
in the word-to-quotient map would otherwise affect both the mathematical
description and its tests. Define
\[
 L=\begin{pmatrix}1&1\\0&1\end{pmatrix},\qquad
 R=\begin{pmatrix}1&0\\1&1\end{pmatrix},\qquad
 \rho(a)=L,\quad \rho(b)=R^{-1}.
\]
The braid relation holds. This is reduced Burau at \(t=-1\).
It is onto \(\SL_2(\Z)\), with kernel
\(\langle z^2\rangle\); \(e(z^2)=12\).

\begin{lemma}[Small-trace conjugacy]
\label{lem:trace}
There are two conjugacy classes in \(\SL_2(\Z)\) with trace \(1\),
represented by \(\rho(ab)\) and its inverse. There is one conjugacy class
with trace \(3\), represented by \(LR=\rho(ab^{-1})\).
\end{lemma}
\begin{proof}
For trace \(1\), Cayley--Hamilton gives \(M^3=-I\).
Its image in \(\PSL_2(\Z)=C_2*C_3\) has order three.
A finite-order element of a free product is conjugate into a factor:
otherwise its cyclically reduced word has length at least two and all its
positive powers remain reduced and grow. The two order-three classes are
therefore the two nonidentity elements of \(C_3\).
Choosing the lift with trace \(1\) fixes the sign and yields the two stated
classes.

For trace \(3\), use
\[
 S=\begin{pmatrix}0&-1\\1&0\end{pmatrix},\qquad U=SL.
\]
In the projective quotient, \(S\) has order two and \(U\) order three.
We have \(SU=-L\) and \(SU^{-1}=R\). A cyclically reduced
infinite-order word is an alternating product of \(S\) and \(U^{\pm1}\).
Consequently a lift is conjugate, up to sign, to a positive word in \(L,R\).
A word in only one letter has trace \(2\). If both occur, rotate the
cyclic word to begin with \(L^rR^s\), \(r,s\ge1\).
This prefix has trace \(2+rs\) and strictly positive entries.
For a strictly positive matrix \(X\),
\[
 \tr(XL)=\tr(X)+X_{21},\qquad
 \tr(XR)=\tr(X)+X_{12}.
\]
Every additional letter strictly increases the trace. Trace \(3\)
therefore forces exactly \(LR\), up to cyclic rotation.
The negative lift has negative trace, so trace \(3\) selects the positive
one. A projective conjugator lifts to an integral determinant-one
conjugator, completing the proof.
\end{proof}

\begin{theorem}[Exact trace--exponent criterion]
\label{thm:matrix}
For a validated three-braid knot,
\begin{equation}
\label{eq:matrixcriterion}
 \widehat\beta\text{ is the unknot}
 \quad\Longleftrightarrow\quad
 (e(\beta),\tr\rho(\beta))
 \in\{(2,1),(-2,1),(0,3)\}.
\end{equation}
\end{theorem}
\begin{proof}
Necessity follows by evaluating the three classes of
Theorem~\ref{thm:classical}. For sufficiency, Lemma~\ref{lem:trace}
gives a matrix conjugacy to the appropriate representative.
For the trace-one alternatives, the two classes have exponent-sum
residues \(2\) and \(-2\) modulo \(12\); the residue is well-defined
because the kernel has exponent \(12\).
The specified exponent therefore selects the correct class.
Lift the matrix conjugator to \(B_3\). The residual difference is a power
of \(z^2\), and equality of exponent sums forces that power to be zero.
\end{proof}

\subsection{The determinant interpretation and a useful trap}

For a three-braid knot, the Alexander--Burau determinant identity gives
\begin{equation}
\label{eq:det}
 \det(\widehat\beta)
   =|\det(I-\rho(\beta))|
   =|2-\tr\rho(\beta)|.
\end{equation}
This provides a further independent test against the existing exact
Alexander implementation. The trace--exponent specialization has historical
precedents in determinant-one closed-three-braid classifications
\cite{birmanhilden,datta2022}; it is not a new classification claim.

The exponent is indispensable. The torus knot represented by
\((ab)^5\) has trace \(1\) and determinant \(1\), but exponent sum \(10\).
More generally, the words \((ab)^{6k+1}\) have the same quotient class
and determinant as \(ab\), while their exponent sums are \(12k+2\).
A classifier that keeps only the quotient or determinant would give
incorrect positive answers on this family.

The figure-eight word \((ab^{-1})^2\) has exponent zero, matrix
\((LR)^2\) of trace \(7\), and determinant \(5\).
Its quotient cyclic word has eight syllables, so both implemented
criteria reject it.

\subsection{Bit complexity and serialization}

Writing a current matrix as
\(\begin{psmallmatrix}u&v\\w&t\end{psmallmatrix}\), right multiplication
by the four generators changes one column by adding or subtracting the
other. For example, multiplication by \(a\) gives
\[
 (u,v,w,t)\longmapsto(u,u+v,w,w+t),
\]
and multiplication by \(b\) gives
\[
 (u,v,w,t)\longmapsto(u-v,v,w-t,t).
\]
After \(j\) letters, the absolute entry sum is at most \(4\cdot2^j\).
Every entry therefore has \(O(j+1)\) bits. Summing the costs of the
constant number of additions per letter yields \(O(\ell^2)\)
elementary bit operations and \(O(\ell)\) bits of matrix storage.
The input word and witness require additional storage if retained.

The implementation encodes large matrix entries in signed hexadecimal
when producing JSON. This avoids decimal-conversion limits and avoids
mistaking a serialization failure for a mathematical failure.
The free-product backend is the default; the matrix backend supplies
independent verification and useful arithmetic witnesses.

\section{Extending the useful input class by certified preprocessing}
\label{sec:markov}

\subsection{A linear braid-writhe obstruction}

Bennequin's theorem implies that an unknotted closed \(m\)-braid has
\(|e(\beta)|<m\), equivalently \(|e(\beta)|\le m-1\)
\cite[pp.~100--101]{bennequin}. Hence
\[
              |e(\beta)|>m-1\quad\Longrightarrow\quad
              \widehat\beta\text{ is nontrivial}.
\]
This is an exact one-sided test for every strand number.
It never promotes equality to an unknot verdict. It is particularly cheap
on long positive or negative braid words and can avoid the general
Alexander matrix entirely.

For one strand, the valid empty word is trivial. For two strands,
\(B_2\cong\Z\), and a knot closure is the unknot exactly when
the exponent sum is \(1\) or \(-1\). Together with
Theorem~\ref{thm:symbolic}, this completes the direct route for at most
three strands.

\subsection{Endpoint singleton destabilizations}

\begin{lemma}[A checked Markov step]
\label{lem:singleton}
Suppose a word on \(m>3\) strands contains exactly one occurrence of a
generator with absolute index \(m-1\). Delete that occurrence after a
cyclic rotation placing it last. The remaining word is an
\((m-1)\)-braid with isotopic closure. The corresponding statement holds
for an absolute-index-one singleton, after deleting the leftmost strand
and decreasing all remaining absolute generator indices by one.
\end{lemma}
\begin{proof}
If the original word is \(u\sigma_{m-1}^{\varepsilon}v\), the absence of
other endpoint generators gives \(u,v\in B_{m-1}\).
Cyclic conjugation gives \(vu\sigma_{m-1}^{\varepsilon}\).
A Markov destabilization removes the last crossing and last strand.
For the left endpoint, conjugation by the half twist reverses the strand
indices, reducing to the right-endpoint case; reversing the remaining
indices back produces the stated renumbering.
Both crossing signs are allowed.
\end{proof}

The routine first freely reduces and cyclically freely reduces the Artin
word. These operations are braid equalities or conjugations.
It then searches only for the two sufficient endpoint conditions.
After each successful destabilization it repeats, stopping at three
strands or when no such step is available.

The certificate records the current strand count and reduced word length,
the chosen endpoint, the occurrence position and signed generator, and
the final word. Its verifier recomputes free reductions, checks that the
chosen endpoint really occurs once, performs the chosen rotation and
deletion, and checks the final result. The verifier does not run the
search that selected the endpoint.

\begin{proposition}[Certified descent bound]
\label{prop:descent}
For an explicit valid word of length \(\ell\) on \(m\) strands, the
implemented singleton descent uses \(O(m\ell)\) word operations,
never increases length, and performs at most \(\max\{m-3,0\}\) destabilizations.
Its certificate can be replayed within the same bound. If it reaches
three or fewer strands, the subsequent classifier gives a complete
decision for the original closure.
\end{proposition}
\begin{proof}
Each pass performs linear free/cyclic reduction and counts endpoint
occurrences. Each successful pass reduces the strand count by one.
There is at most one final unsuccessful pass. Rotation and renumbering
cost \(O(\ell)\); all later words have length at most the original.
The same bound holds for replay, which checks rather than searches for
each chosen occurrence. Soundness follows by composing
Lemma~\ref{lem:singleton} and the free reductions, and then applying the
complete low-strand criterion.
\end{proof}

Since \(m\le\ell+1\) for a one-component input, the routine has a
polynomial worst-case bound. This extends the certified family beyond
inputs initially declared to have three strands, without claiming an
algorithm for minimal braid index.

\subsection{A necessary negative control: Morton's example}

Morton's four-braid unknot \cite{morton}, in the word convention displayed
by Birman--Brendle \cite[Example 3.1]{birmansurvey}, is
\[
 X=\sigma_3^{-2}\sigma_2\sigma_3^{-1}\sigma_2
      \sigma_1^3\sigma_2^{-1}\sigma_1\sigma_2^{-1}.
\]
Its explicit word is
\[
 [-3,-3,2,-3,2,1,1,1,-2,1,-2].
\]
The exponent sum is \(1\), so the writhe test is inconclusive.
The left and right endpoint generators occur four and three times
respectively. The singleton routine therefore stops immediately.
The new preprocessing returns no verdict and the existing general
recognizer supplies the unknot answer. This example is in the tests.

More generally, a failure to descend says nothing about whether a closure
is knotted or whether a different representative has fewer strands.
Exchange moves and other transformations can be necessary.

\subsection{A closure certificate is not a local tangle replacement}

An important boundary of applicability is easy to miss. A three-braid
whose \emph{closure} is an unknot need not be the identity braid.
For example \(ab\) has unknot closure and exponent two, so it is not the
identity. Replacing occurrences of it by the identity inside another
braid can change the resulting knot.

The new classifier is used on a whole verified closed braid, or after
whole-closure-preserving Markov steps. It cannot be used to erase an
arbitrary three-strand subword of a larger tangle. Likewise, the existing
connected-sum shortcut is justified by an actual splitting sphere, not
merely by recognizing an auxiliary closure.

\section{Integration and exact semantics}
\label{sec:implementation}

The parser previously converted braid input to a normalized PD code and
discarded the word. The updated parser retains immutable validated braid
provenance. PD equality and hashing continue to refer to the same diagram
data. The new decision route is invoked before geometric simplification
discards that provenance.

The implementation validates provenance rather than accepting a user
assertion that an unrelated PD has a three-braid representation.
Serialization, restoration, and mirror handling are tested.
The direct word API is also available for large explicit words when no
PD construction is needed.

The resulting order of work is:
\begin{enumerate}
\item Validate a one-component input and any retained braid provenance.
\item When braid provenance is available, apply the writhe obstruction
and the complete route for at most three strands.
\item For unresolved higher-strand inputs, first run the existing cheap
Reidemeister I/II and descending checks, then optional certified endpoint
descent on the original validated source word. The original closure and
the simplified diagram are isotopic, so this certificate remains valid.
\item If these do not finish, use the existing polynomial, factorization,
Reidemeister III, and Khovanov pipeline.
\end{enumerate}
Delaying higher-strand descent avoids repeated word passes on easy kink
chains already handled by the cheap geometric checks. The raw word API
still offers descent directly, since it does not construct a PD diagram.
Each added early decision is justified independently of the fallback.
The fallback is unchanged mathematically. Its complete decision theorem
is reduced Khovanov unknot detection \cite{kronheimer}: rank one over
\(\F\) occurs exactly for the unknot. Indeed rational reduced rank is
at most the rank over \(\F\) by universal coefficients, and rational
rank one detects the unknot; the unknot itself has rank one over both fields.

A resource limit remains a third outcome, \code{UNKNOWN}.
It is not a negative answer. Backend options allow the new route to be
disabled when testing or benchmarking the general scanner.
Resource-limit regression tests that intentionally force that scanner
explicitly turn off the new route; otherwise a valid early certificate
would correctly finish before the intended limit.

\subsection{Code delivered}

\begin{center}
\begin{tabularx}{\textwidth}{@{}p{0.36\textwidth}X@{}}
\toprule
Component & Purpose\\
\midrule
Braid backend module & Direct symbolic and matrix decisions, writhe witness,
 component checks and statistics\\
\code{braid_reduction.py} & Elementary Markov descent and independent
 certificate replay\\
\code{residue.py} & Non-mutating scalar-residue multiplicity diagnostic\\
Parser and recognizer integration & Retained provenance, options,
 serialization and exact fallback behavior\\
Tests and experiment scripts & Exhaustive comparisons, independent oracles,
 paired timings and intrinsic-profile checks\\
\bottomrule
\end{tabularx}
\end{center}

The mathematical distinction between decision and full invariant output
is retained in the API. A braid certificate proves a verdict; it is not a
request to fabricate uncomputed Khovanov groups.
The existing full-rank APIs continue to mean full-rank computations.

\section{Intrinsic multiplicities of a partial scanning complex}
\label{sec:residue}

The second main result addresses the general scanner rather than only
braid input. We work over \(\F\), as the implementation does.
Fix one prefix, its boundary, and its matching objects.
A complex is a finite direct sum of matchings in each homological degree,
with differential entries in the dotted cobordism algebra.

For a boundary with \(2b\) points and two matchings \(a,c\), let
\(k(a,c)\) be the number of circles in their alternating union.
A morphism has a basis indexed by dot subsets on these circles.
In particular
\[
             \End(a)\cong
       \F[x_1,\ldots,x_b]/(x_1^2,\ldots,x_b^2).
\]
Its units are exactly the polynomials with constant coefficient one.

\subsection{The scalar residue functor}

Define \(q\) on a morphism by discarding it when the two matching objects
differ, and by taking its undotted coefficient when they agree.
The target is a category with one independent scalar identity for each
matching.

\begin{lemma}[Residue respects composition]
\label{lem:residuefunctor}
For the matching cobordism algebra used by the scanner, \(q\) is additive,
preserves identities, and satisfies \(q(gf)=q(g)q(f)\).
\end{lemma}
\begin{proof}
Additivity and identities are immediate. A composition with distinct
initial and final matchings is discarded on both sides.
For \(a\to a\to a\), composition is multiplication in the displayed
square-zero polynomial algebra, and constant terms multiply.

It remains to consider \(a\to c\to a\) with \(c\ne a\).
Then \(k(a,c)<b\): having \(b\) alternating-union circles forces every
arc to agree. In the gluing surface, every connected component contains
an input circle of the left morphism. Thus there are at most \(k(a,c)\)
components, whereas the output consists of \(b\) circles.
An undotted constant output can arise only when each surface component
has exactly one output circle and no dot: a genus-zero component with
two or more output circles contributes at least one dot, and positive-genus components vanish over \(\F\).
A component with no output cannot increase the number of components
available for output. A constant output would require at least \(b\)
components, contradicting \(k(a,c)<b\).
Therefore its undotted coefficient vanishes.
\end{proof}

The proof uses the algebra of the fixed prefix and does not require that
all boundary points lie on one disk boundary.
For disk tangles there is also the shorter positive-grading proof in
Section~\ref{sec:batch}.

Let \(N_{h,a}\) be the number of copies of matching \(a\) in cube degree
\(h\), and let \(D_{h,a}\) be the scalar differential matrix
from \((h,a)\) to \((h+1,a)\) after applying \(q\).
By the lemma, \(D_{h+1,a}D_{h,a}=0\). Set
\begin{equation}
\label{eq:mu}
 \mu_{h,a}(C)
 =N_{h,a}-\rank D_{h,a}-\rank D_{h-1,a},\qquad
 \mu(C)=\sum_{h,a}\mu_{h,a}(C).
\end{equation}

\begin{theorem}[Canonical post-cancellation multiplicities]
\label{thm:canonical}
Any exhaustive sequence of invertible-entry cancellations in \(C\)
leaves exactly \(\mu_{h,a}(C)\) copies of matching \(a\) in degree \(h\).
These numbers depend only on the chain-homotopy type with its fixed
homological normalization.
\end{theorem}
\begin{proof}
Categorical Gaussian cancellation is a chain-homotopy equivalence
\cite{barnatan}. Applying an additive functor preserves chain maps,
homotopies, and their equations, so applying \(q\) to a cancellation
preserves the homology of the scalar quotient complex.

Once no invertible entry remains, every same-matching entry has constant
term zero and every different-matching entry is killed by \(q\).
The residue differential of the remaining complex is therefore zero.
Its number of copies in each \((h,a)\) is its residue homology dimension.
That dimension agrees with the original complex and is exactly
\eqref{eq:mu}. The same argument proves invariance under any
chain-homotopy equivalence with the stated normalization.
\end{proof}

Khovanov's minimal-core results supply the general algebraic background
\cite[Section 5.1]{khovanovtangles}.
Equation~\eqref{eq:mu} is the concrete diagnostic appropriate to this code.

\subsection{What pivot optimization can change}

The theorem gives a sharp interpretation of pivot experiments.
At the \emph{same completed prefix}, every exhaustive pivot strategy
has the same remaining object multiplicities. A better strategy can
still reduce:
\begin{itemize}
\item the number of compositions and the size of intermediate differential
fill-in;
\item temporary storage during a cancellation stage;
\item bookkeeping and cache overhead.
\end{itemize}
It cannot reduce the final minimal object count at that prefix.

Changing scan order changes the partial tangles, and therefore can change
their intrinsic multiplicities. Delaying cancellation across several
crossings also changes the sequence of completed prefixes being
materialized. These remain meaningful directions; the theorem does not
say that every ordering has the same complexity.

The code reports cube degree. Reidemeister changes can shift the usual
normalized homological degree, so degree-by-degree comparisons across
different presentations must align conventions. Our pivot comparisons
use identical prefixes and require no such shift.

\subsection{Implemented diagnostic and its limitation}

The diagnostic builds scalar blocks, computes their \(\F\)-ranks with
bit-packed elimination, and returns the nonzero \(\mu_{h,a}\).
It accepts the optimized scanner and the set- and bit-based generic
scanners. It optionally checks the scalar equation \(d^2=0\), verifies
degree consistency, and does not mutate the complex.

It is deliberately not an extra default pass before every elimination:
that would duplicate work unless the profile is being used for an audit
or a different batching algorithm. Its purpose in this bundle is to
verify an invariant, make measurements interpretable, and provide a
basis for subsequent algorithms.

Crucially, the multiplicity list alone is not a replacement for the
partial complex. Positive-degree differential entries survive in the
minimal complex. Later crossings can interact with them. Dropping those
entries and retaining only the list would not be an exact Khovanov
algorithm.

\section{Finite batch cancellation through the radical}
\label{sec:batch}

This section develops a concrete next algebraic primitive, with a proof
and cost bound. It is not described as an implemented production speedup.

Assume first that the prefix is represented as a disk tangle with
\(2b\) boundary points. The standard arc algebra is
\[
 A_b=\bigoplus_{a,c\in\mathcal M_b}\Hom(a,c),
\]
where \(\mathcal M_b\) is the set of crossingless disk matchings.
A basis element with \(d\) dots has intrinsic algebra degree
\begin{equation}
\label{eq:arcdegree}
              \deg_{\rm alg}=b-k(a,c)+2d.
\end{equation}
This is the positive grading from the arc-algebra construction
\cite[Section 2.4]{khovanovtangles}; it is distinct from the
homological degree and from presentation shifts in quantum degree.

\begin{lemma}[A bounded radical filtration]
\label{lem:radical}
The algebra \(A_b\) is graded in degrees \(0,\ldots,2b\).
Its degree-zero part is \(\prod_{a\in\mathcal M_b}\F e_a\).
Its positive-degree ideal \(J_b\) satisfies \(J_b^{2b+1}=0\);
projection onto degree zero is the residue functor.
\end{lemma}
\begin{proof}
Since \(k\le b\) and \(0\le d\le k\), \eqref{eq:arcdegree} lies between
zero and \(2b\). It is zero only for \(a=c\) and the undotted identity.
Composition preserves degree. This can also be checked from the
Frobenius operations: every split adds one dot in each nonzero summand,
whereas a merge adds none; counting the \(b\) saddles in a composition
gives additivity of \eqref{eq:arcdegree}.
Thus positive degree is an ideal, and a product of \(2b+1\) positive
terms must vanish. The quotient is the displayed product of fields.
Over \(\F\), this ideal is also the Jacobson radical.
\end{proof}

\Needspace{16\baselineskip}
\begin{theorem}[Finite simultaneous reduction]
\label{thm:hpl}
Let \(d=d_0+\delta\), where \(d_0\) consists of the scalar
equal-matching coefficients of a partial complex. Choose a scalar
contraction of \((C,d_0)\) onto its zero-differential homology object \(H\):
\[
 pi=1_H,\quad d_0h+hd_0=1_C+ip,\quad
 h^2=hi=ph=0,\quad d_0i=pd_0=0.
\]
All entries of \(i,p,h\) can be scalar matching identities.
Define
\[
 {\cal T}
   =\delta\sum_{j=0}^{2b}(h\delta)^j,\qquad
                    d_H=p{\cal T}i.
\]
Then \((H,d_H)\) is chain-homotopy equivalent to \((C,d)\),
has exactly the multiplicities \(\mu_{h,a}(C)\), and is fully
unit-reduced.
\end{theorem}
\begin{proof}
The residue functor implies \(d_0^2=0\).
Scalar Gaussian elimination in each matching block supplies the stated
special contraction. The entries of \(\delta\) lie in \(J_b\), while
\(h\) has algebra degree zero. Hence
\((h\delta)^{2b+1}=0\), so the displayed sum is the exact finite inverse
of \(1+h\delta\), multiplied on the left by \(\delta\).
There is no convergence or numerical approximation issue.

For completeness, put
\[
 i'=i+h{\cal T}i,\qquad p'=p+p{\cal T}h,\qquad
 h'=h+h{\cal T}h.
\]
Over \(\F\), the finite perturbation identities are
\[
 {\cal T}=\delta+\delta h{\cal T}
         =\delta+{\cal T}h\delta,\qquad
 d_0{\cal T}+{\cal T}d_0={\cal T}ip{\cal T}.
\]
The second follows by inserting \(d_0\delta+\delta d_0=\delta^2\)
and \(d_0h+hd_0=1+ip\) in the finite series.
They imply
\[
 d_H^2=0,\quad di'=i'd_H,\quad p'd=d_Hp',\quad
 p'i'=1_H,\quad dh'+h'd=1_C+i'p'.
\]
These are the required chain-homotopy equations.
They are the homological perturbation formulas specialized to the
nilpotent ideal \cite{crainic}.
Every term of \(d_H\) contains a \(\delta\), so its entries have positive
algebra degree. It therefore has no cancellable unit.
\end{proof}

For \(b=0\), the radical is zero and this reduces to ordinary homology.
For \(b>0\), computing only \(H\) is insufficient: the transferred
\(d_H\) must be retained. The nilpotence index bounds the number of
matrix products, not the size of their coefficient matrices.

\subsection{A conservative explicit arithmetic bound}

Let \(M\) be the number of objects before a stage is reduced and
\(D=\dim_{\F}A_b\). If \(\omega\) is an admissible exponent for matrix
multiplication over \(\F\), a direct dense realization of
Theorem~\ref{thm:hpl} uses
\[
                  O\bigl((b+1)D^3M^\omega\bigr)
\]
field operations, in addition to building the finite multiplication
table. Expand entries in an algebra basis: one algebra-valued matrix
product is at most \(D^3\) coefficient-matrix products. There are
\(O(b+1)\) such products in the formula. Scalar rank decomposition
supplies a contraction within \(O(M^\omega)\) field operations.

For the disk algebra,
\[
 D=\sum_{a,c}2^{k(a,c)}\le\Cat(b)^2\,2^b=2^{O(b)}.
\]
Thus the dense bound is \(2^{O(b)}M^\omega\).
This can be poor on sparse small instances; it is an asymptotic
implementation target, not a claim that dense multiplication will beat
the current sparse scanner in the supplied measurements.

The disk hypothesis is real. For an arbitrary multi-boundary prefix,
one must specify its actual matching algebra or prove a disk
description. A safe count of abstract matchings on \(2b\) points is
\((2b-1)!!\), not automatically \(\Cat(b)\). Planarity of the full knot
diagram by itself does not supply the missing boundary description.

\Needspace{9\baselineskip}
\section{What a width--multiplicity bound really proves}
\label{sec:complexity}

Let a specified crossing order have at most \(2b\) active boundary
points, and let
\[
             \mu=\max\{1,\mu(C_i):0\le i\le n\}
\]
be its largest intrinsic minimal multiplicity.
After a crossing is added to a reduced prefix, there are two smoothings,
and each can create at most two closed circles. Delooping therefore
creates at most \(8\mu\) objects before the next reduction.

\begin{theorem}[A safe scalar profile bound]
\label{thm:profile}
There is an exact scalar-cancellation realization using the same matching
representation whose cost, for the specified scan, is
\[
          \softO\bigl(n\,2^{O(b)}\mu^3\bigr).
\]
In particular,
\[
                   b+\log(\mu+1)=O(\log^2(n+2))
\]
is sufficient for a quasi-polynomial bound of the form
\(n^{O(\log n)}\).
\end{theorem}
\begin{proof}
At one stage there are \(M=O(\mu)\) objects. Locate a unit by scanning
at most \(M^2\) entries. A cancellation removes two objects and changes
at most \(M^2\) differential entries. There are at most \(M/2\)
cancellations. Each morphism has at most \(2^b\) dot monomials and
composition can be evaluated by an explicit \(2^{O(b)}\operatorname{poly}(b)\)
procedure. The entry scans, compositions and bit-packed additions
therefore cost \(2^{O(b)}M^3\), up to logarithmic bookkeeping.
Transfer at a crossing costs no more at this scale.
Sum over \(n\) stages.
Taking logarithms of the bound proves the final assertion.
\end{proof}

The theorem describes a definite available scalar algorithm.
It is not an unproved assertion that every stale-priority queue or
racing variant in the current code has exactly this overhead bound.
The measured intrinsic profile remains meaningful for those variants.

For disk frontiers, the batch formula gives
\(n2^{O(b)}\mu^\omega\), with the same conceptual limitation.
In both cases, the problem is establishing or exploiting a useful
bound on \(\mu\). A measured small profile for a benchmark is not a proof
of such a bound on all knots.

\subsection{An explicit fixed-width obstruction}

For \(r\ge1\), consider the alternating three-braid word
\[
                  \beta_r=(ab^{-1})^r.
\]
For \(3\nmid r\), its closure is a knot. The crossing order inherited
from the braid word has at most six frontier endpoints: at most three at
the moving level and three associated with the closure.
Set \(\lambda_\pm=(3\pm\sqrt5)/2\).
Since \(LR\) has eigenvalues \(\lambda_\pm\),
\[
 \det(\widehat{\beta_r})
 =\tr((LR)^r)-2
 =\lambda_+^r+\lambda_-^r-2.
\]
The reduced Khovanov Euler characteristic is the normalized Jones
polynomial, with the usual square-root change of variable. Its graded
Euler characteristic is unchanged by passing to \(\F\). In a convention
where the Jones variable is \(t=q^2\), evaluate at \(q=i\); then
\(|V_K(-1)|=\det(K)\), and every monomial has absolute value one.
The triangle inequality therefore bounds the determinant by the total
reduced Betti number, giving
\[
 \dim_{\F}\widetilde{\Kh}(\widehat{\beta_r})
       \ \ge\ \det(\widehat{\beta_r}).
\]
Consequently the size of an explicit minimal complex grows exponentially
in \(r\), despite the fixed number of braid strands.
This elementary family illustrates the representation obstruction
analyzed more broadly in \cite{kelomaki2026}.

This is a statement about materializing the homology or its minimal
complex. The baseline's determinant filter already rejects this family
cheaply. We do not use a deliberately disabled filter to claim an
end-to-end speedup of the new pipeline. The relevant lesson is that an
explicit-basis fallback cannot acquire a general polynomial width bound
solely from improved cancellation bookkeeping.

\subsection{Decision, output size, and compressed multiplicity}

An exponentially large rank can be stored as an integer with linearly
many bits. An explicit basis with that many elements cannot.
The existing connected-sum shortcut already uses this distinction.
The three-braid classifier takes it further by computing only a
complete decision invariant for its class.

For general tangles the difficult information lies in the differential,
not only in the counts of objects. A useful compressed representation
must support exact gluing, composition, and rank or equivalence tests
without expanding the multiplicities. The residue theorem says what
survives scalar cancellation; it does not make those operations free.

\section{Experiments and verification}
\label{sec:experiments}

All empirical claims below concern completed runs on the supplied
implementation and pinned baseline. The experiment artifacts record
the interpreter, platform, options, exact inputs or seeds, operation
counts, and raw paired timings. Wall-clock results are evidence about
these workloads, not proofs of complexity.

\subsection{Independent correctness checks}

The pinned baseline passes 41 unit tests; the delivered implementation
passes 59. These test counts include the following exhaustive comparisons,
not merely a handful of selected examples.

The new three-braid tests exhaust all \(46{,}376\) one-component words
of lengths two, four, six and eight over
\(\{a,a^{-1},b,b^{-1}\}\), comparing the free-product and matrix criteria.
For all \(2{,}856\) such words through length six, an independent dense
reduced-Khovanov cube is used as another decision oracle.

Additional checks compare the determinant formula against the existing
exact Alexander polynomial on 120 seeded words; test braid
relations and conjugation in 200 seeded contexts; stress long words and large matrix
serialization; and exercise provenance, mirrors, invalid input,
resource limits and fallback behavior. These compare different
mathematical computations rather than only duplicate implementations of
one predicate.

The Markov tests compare unreduced input and reduced output through
the existing Khovanov engine for both endpoint choices and both crossing
signs. They test repeated stabilizations under conjugation and reject
tampered certificates. Morton's example is a mandatory inconclusive
preprocessing case.

For multiplicities, every tested crossing prefix is transferred without
elimination, its scalar-residue Betti numbers are computed, and then the
actual categorical cancellations run. The predicted and observed
matching-by-degree profiles must agree, including across the optimized,
generic bit-based, and generic set-based engines with distinct pivot
rules. An additional uncancelled closed complex is checked against the
existing full linear-homology routine.

\subsection{Paired performance measurements}

The benchmark separates direct word processing from public pipeline
calls on parsed diagrams. It also separates the baseline default
pipeline from a named ablation that disables Reidemeister III search.
The latter is useful for exposing a scanner workload, but is not
substituted for the default baseline comparison.

% BENCHMARK_TABLE_BEGIN
The final run used CPython 3.12.14 on x86\_64 Linux, with 15 paired
rounds per case and a rotating three-arm order: baseline, new, and a
second baseline call. Each implementation had one warm-up. The timed
region includes fresh PD construction and recognition; it excludes
process startup, imports, and JSON decoding. All published samples
finished with the expected exact verdict.

The comparison rule uses an order-statistic interval with at least 95\%
coverage for the median new/old ratio under independent sampling. It
labels a case faster only when the interval lies below one and the
median ratio is below the first quartile of the paired baseline-control
ratios. The slower rule is symmetric. These are exploratory workload
comparisons with a local A/A noise control; they are not simultaneous
confidence guarantees across the whole corpus.

Twenty-one of 28 cases satisfy the faster rule, seven support no timing
claim, and none satisfy the slower rule. This is a statement about these
completed measurements, not a proof that regressions are impossible.
The default-pipeline speedup ranges are 3.97--7.04 for positive
three-braids, 2.39--3.88 for alternating three-braids, and 2.44--5.29
for the five scrambled three-braid unknots. The larger 10.72--51.05
range belongs to the explicitly labeled R3-disabled ablation.

\begingroup
\small
\setlength{\tabcolsep}{4pt}
\begin{longtable}{@{}p{0.38\textwidth}rrrrl@{}}
\caption{All paired pipeline measurements. Times are separate medians in
milliseconds; speedup is the median of paired old/new ratios, which need
not equal the quotient of the displayed medians. ``None'' means no timing
claim under the stated rule.}\label{tab:paired}\\
\toprule
Case & \(n\) & Old ms & New ms & Speedup & Claim\\
\midrule
\endfirsthead
\toprule
Case & \(n\) & Old ms & New ms & Speedup & Claim\\
\midrule
\endhead
\midrule
\multicolumn{6}{r}{\emph{Continued on next page}}\\
\endfoot
\bottomrule
\endlastfoot
Positive three-braid, \(r=5\) & 10 & 0.351 & 0.087 & 4.09\(\times\) & Faster \\
Positive three-braid, \(r=11\) & 22 & 0.630 & 0.156 & 3.97\(\times\) & Faster \\
Positive three-braid, \(r=23\) & 46 & 1.559 & 0.272 & 4.40\(\times\) & Faster \\
Positive three-braid, \(r=47\) & 94 & 3.966 & 0.575 & 7.04\(\times\) & Faster \\
Positive three-braid, \(r=95\) & 190 & 8.705 & 1.452 & 5.67\(\times\) & Faster \\
Alternating three-braid, \(r=5\) & 10 & 0.210 & 0.085 & 2.39\(\times\) & Faster \\
Alternating three-braid, \(r=11\) & 22 & 0.448 & 0.158 & 2.91\(\times\) & Faster \\
Alternating three-braid, \(r=23\) & 46 & 2.040 & 0.424 & 3.88\(\times\) & Faster \\
Alternating three-braid, \(r=47\) & 94 & 2.752 & 0.871 & 2.52\(\times\) & Faster \\
Alternating three-braid, \(r=95\) & 190 & 7.481 & 2.244 & 2.72\(\times\) & Faster \\
Scrambled unknot, seed 0 & 190 & 6.059 & 2.221 & 2.44\(\times\) & Faster \\
Scrambled unknot, seed 0 (R3 off) & 190 & 17.365 & 1.731 & 10.72\(\times\) & Faster \\
Scrambled unknot, seed 4 & 158 & 6.295 & 1.197 & 5.12\(\times\) & Faster \\
Scrambled unknot, seed 4 (R3 off) & 158 & 31.022 & 1.502 & 21.36\(\times\) & Faster \\
Scrambled unknot, seed 5 & 202 & 10.292 & 2.411 & 4.42\(\times\) & Faster \\
Scrambled unknot, seed 5 (R3 off) & 202 & 80.592 & 1.563 & 51.05\(\times\) & Faster \\
Scrambled unknot, seed 7 & 198 & 6.787 & 1.394 & 4.36\(\times\) & Faster \\
Scrambled unknot, seed 7 (R3 off) & 198 & 58.837 & 2.089 & 27.66\(\times\) & Faster \\
Scrambled unknot, seed 14 & 182 & 6.815 & 1.259 & 5.29\(\times\) & Faster \\
Scrambled unknot, seed 14 (R3 off) & 182 & 19.499 & 1.185 & 14.83\(\times\) & Faster \\
Morton's four-braid unknot & 11 & 2.818 & 2.913 & 1.03\(\times\) & None \\
Stabilized unknot, 16 strands & 15 & 0.120 & 0.133 & 0.93\(\times\) & None \\
Stabilized unknot, 64 strands & 63 & 0.474 & 0.544 & 0.81\(\times\) & None \\
Stabilized unknot, 256 strands & 255 & 2.838 & 2.206 & 1.20\(\times\) & None \\
Conway & 11 & 0.820 & 0.795 & 0.97\(\times\) & None \\
Kinoshita--Terasaka & 11 & 0.932 & 0.869 & 1.04\(\times\) & None \\
Eight-crossing unknot (braid) & 8 & 0.508 & 0.092 & 5.57\(\times\) & Faster \\
Grid unknot & 6 & 0.080 & 0.078 & 1.02\(\times\) & None \\
\end{longtable}
\endgroup

Raw word timings are separate five-repetition experiments. At 128,000
letters the symbolic median is 0.112884 seconds and the matrix median
is 0.626293 seconds; the latter's largest final entry has 88,863 bits.
At 1,024,000 letters the symbolic median is 1.044478 seconds, with exactly
2,048,000 quotient-token operations. At 1,000 letters, however, the matrix
code is faster (0.000258 versus 0.000744 seconds). The asymptotically better
symbolic method does not win every small case. Raw scaling samples do not
use the three-arm comparison rule and are reported as descriptive timings.

% BENCHMARK_TABLE_END

For repeated calls, baseline and new calls are interleaved, and both
orders occur. Ratios are computed from paired observations.
Parsing and fresh-process startup are reported only when included by
the named workload. The raw data, rather than a cross-machine historical
table, are the authority for these ratios.

\subsection{Deterministic complexity evidence}

Operation counts test the linear word-processing bound independently of
timer noise. Long conjugated words and words whose reduced quotient
stays long exercise different stack and cyclic-reduction behavior.
The matrix experiment records the largest final matrix-entry bit length and contrasts
its larger-integer cost with constant-alphabet reduction.

A second experiment records intrinsic multiplicities, boundary size,
cancellation counts and composition counts for identical prefix orders.
It checks that pivot changes can alter work while preserving the
canonical completed-prefix multiplicities. The data do not extrapolate
small observed profiles into a universal bound.

The deterministic profile experiment includes 54 prefixes on six diagrams,
checked in four scanner configurations: 216 pre/post-cancellation checks.
Every predicted multiplicity and every completed profile agrees.
The table compares composition calls in the two generic bit-based engines,
so their counters have the same meaning; optimized identity shortcuts are
not mixed into this comparison.

\begin{center}
\small
\begin{tabular}{@{}lrrrr@{}}
\toprule
Diagram & Crossings & Peak \(\mu\) & Min-fill calls & LIFO calls\\
\midrule
Conway & 11 & 66 & 770 & 1281 \\
Kinoshita--Terasaka & 11 & 66 & 671 & 1277 \\
Figure-eight & 4 & 10 & 15 & 18 \\
Eight-crossing unknot & 8 & 4 & 13 & 24 \\
\(T(3,5)\) & 10 & 14 & 113 & 132 \\
\((ab^{-1})^5\) & 10 & 242 & 1216 & 1707 \\
\bottomrule
\end{tabular}
\end{center}
For example, both Conway runs retain the same maximum of 66 objects,
while their composition counts differ. This is precisely the distinction
proved by Theorem~\ref{thm:canonical}.

\subsection{Limits of the experimental evidence}

The new complete route requires supplied or certified braid structure.
A PD input that happens to represent a three-braid knot is not
automatically routed there. Geometric discovery of the representation
is outside these measurements.

The general algorithm still has an exponential fallback. The new code
does not implement the radical batching formula, a compressed general
Khovanov differential, or the missing geometric hierarchy engine.
No measured speedup is used as evidence that these open tasks have
already been completed.

\section{The geometric route to the quasi-polynomial target}
\label{sec:hierarchy}

The 2026 hierarchy work is important because it identifies mathematical
operations that a full implementation could represent explicitly.
It should be used with the precise complexity distinction made in its
Section 9 \cite{lackenby2026}.

Write \(L\) for the maximum hierarchy length and \(g\) for the maximum
pattern-complexity \(\chi_p(S)\) of any decomposing surface appearing
during the algorithm. Proposition 9.1 supplies an
iteration bound of the form
\[
                         L(g+1)^L.
\]
An iteration count is not a complete bit-complexity analysis.
It must be combined with representation sizes, preparation costs,
the exact cost of each operation, and its change to later states.

\begin{proposition}[A noncircular complexity-transfer contract]
\label{prop:hierarchy}
Suppose a hierarchy implementation on an \(n\)-crossing input has
initialization cost \(P(n)\), at most \(L(n)(g(n)+1)^{L(n)}\)
iterations, and a proved uniform bit cost \(Q(n)\) for an iteration,
including all data updates and any necessary verification.
Then
\[
 T(n)\le P(n)+L(n)(g(n)+1)^{L(n)}Q(n).
\]
In particular, if \(P(n),Q(n)\le n^{O(\log n)}\),
\(g(n)\le n^{O(1)}\), and \(L(n)=O(\log n)\), then
\(T(n)=n^{O(\log n)}\).
\end{proposition}
\begin{proof}
Sum the uniform cost over the proved number of iterations and add
initialization. Under the stated bounds,
\[
 \log\!\bigl(L(g+1)^LQ\bigr)
    =O(\log\log n)+O((\log n)^2)+O((\log n)^2).
\]
Exponentiating gives the claim.
\end{proof}

Here \(Q(n)\) must be proved using the evolving encoded state, not
defined to be quasi-polynomial by assumption without an encoding
argument. If \(L=O(\log^2 n)\) and \(g\) remains polynomial, this
particular product bound yields \(n^{O(\log^2 n)}\), not the sharper
\(n^{O(\log n)}\) target. Both are often called quasi-polynomial;
the exponents must be kept explicit.

The missing project obligations are concrete:
\begin{enumerate}
\item An encoded handle structure, boundary pattern, and hierarchical
multi-surface with clear provenance and bit-size bounds.
\item Compressed cutting and gluing that do not expand large normal
coordinates into that many physical pieces.
\item Constructive simplification, weak reduction, and the required
geometric region operations, with exact topological contracts.
\item A decreasing potential and a depth or aggregate-work estimate
strong enough for the target exponent.
\item A complete cost analysis for all intermediate state changes.
\end{enumerate}

A polynomially checkable certificate is valuable but does not provide
a polynomial-time certificate constructor. Likewise, an interval
pairing or normal-coordinate primitive does not by itself supply every
hierarchy transformation. These distinctions explain what a next
implementation must prove rather than merely promise.

\section{Research questions and proposed next work}
\label{sec:future}

The following questions are chosen for their effect on the actual
remaining costs. Each includes a concrete success criterion.

\subsection{Recover certified small-braid descriptions from PD inputs}

Can a practical algorithm construct a three-braid representation, or a
sequence of whole-closure-preserving moves leading to one, for a useful
class of PD or grid inputs? The output should be a replayable certificate
with a proved relation between certificate size and original input size.

A useful first result would recognize a nontrivial structural class in
polynomial time, then invoke the backend proved here. Merely noting that
a knot has braid index three does not provide the representation required
by the algorithm.

\subsection{Go beyond singleton endpoint reductions}

Implement certified exchange moves, bounded bands of braid relations,
or a controlled low-strand conjugacy search before attempting
destabilization. Morton's four-braid is a small mandatory example, but
success on it alone is not a general theorem.

A meaningful complexity result should bound total explored states or
supply an independently decreasing potential. The search should retain
its incomplete outcome and hand unhandled cases to the exact fallback.

\subsection{A complete structural decision route for four-braids}

The three-braid quotient has an exceptionally simple central extension.
There is no justified instruction to reuse the same predicate on four
strands. Can four-braid unknot recognition be reduced to finitely many
structural normal forms or to a compressed invariant that is complete
for the trivial knot?

This is more focused than computing every invariant, and complements
the fixed-strand Khovanov-complexity questions in
\cite{kelomaki2026}. A first milestone is a complete subclass with
explicit encodings and a cost proof, together with counterexamples to
overgeneralized trace tests.

\subsection{Use intrinsic multiplicity to choose scan orders}

A frontier-width objective ignores the invariant measured by
\(\mu(C_i)\). Can inexpensive algebraic probes predict that invariant
well enough to select better next crossings? Can one prove approximation
or parameterized guarantees for an order minimizing
\(\max_i\mu(C_i)\), perhaps jointly with boundary size?

The residue diagnostic can label training or benchmark instances,
but computing the exact profile after building the full stage does
not itself give a cheap predictor. Tests should compare actual
composition counts and peak memory, not only a surrogate score.

\subsection{Implement and compare radical batching}

Build the scalar contraction and transferred positive-degree
differential of Theorem~\ref{thm:hpl}. Verify the full
chain-homotopy equations on small real tangle complexes, then compare
sparse, block, and dense versions on stages where scalar elimination
dominates.

The success criterion is an actual reduction in relevant work or
memory, with the original differential information preserved.
The nilpotence bound alone does not guarantee a practical improvement.

\subsection{Compress differential operators, not only ranks}

For repeated braid blocks or repeated tangle patterns, can the minimal
differential be represented by structured matrices, tensor networks,
or recurrences whose composition and final rank are efficiently
computable? A correct representation must survive gluing and
cancellation without expanding an exponential basis.

The first decisive theorem would establish closure of a useful
representation class under those operations with controlled encoded
size. Reporting exponentially large ranks in binary is only the
starting point.

\subsection{Resolve the multi-boundary matching model explicitly}

Specify the boundary topology of arbitrary scanner prefixes and prove
the appropriate state-space bound for that topology.
Where a Catalan estimate is desired, give a disk-frontier certificate
or a transformation whose additional width is accounted for.

This question affects both bracket and categorical cost statements.
Using the conservative matching count is safe, but understanding
which smaller algebra is actually available could give a meaningful
parameter improvement.

\subsection{Prove a useful multiplicity-profile theorem}

Identify a family with
\(b+\log(\mu+1)=O(\log^2 n)\) along an efficiently constructible order.
The theorem should not follow merely from a small number of strands:
the explicit homological output obstruction shows that this is
insufficient.

A more realistic class might combine a decomposition certificate,
controlled local blocks, and a rule preserving compressed differential
structure. Any decomposition must carry a topological gluing proof.

\subsection{Implement one complete geometric hierarchy operation}

Select a single currently missing transformation, specify its encoded
input and output, prove topology preservation, and bound its bit cost
and output size. Compressed cutting with provenance, or a precise
weak-reduction routine, would be a substantial milestone.

Completing one such primitive is preferable to adding another
high-level hierarchy driver whose expensive subroutines remain
unspecified. The subsequent theorem should compose its bound with
the potential and depth analysis of the whole hierarchy.

\subsection{Strengthen certificates and formalize the small trusted core}

The free-product reduction, central-lift lemma, accepted conjugacy
types, matrix products, and endpoint certificate replay form a compact
verification target. A formalization could establish the algebraic
parts independently of the much larger general scanner.

The topological input should remain a clearly named theorem:
formalizing the stack algorithm does not by itself formalize the
closed-three-braid classification. Likewise, a proof checker should
distinguish a word certificate, a diagram-equivalence certificate, and
a whole-closure Markov certificate.

\section{Reproduction and integration}
\label{sec:reproduce}

The archive contains the updated Python source, tests, a code patch
against the pinned baseline, this article and its source, raw results,
experiment scripts, provenance, and checksums. No third-party lecture
slides or source papers are redistributed.

Use the included instructions for the exact commands and backend
options. The principal checks are:
\begin{lstlisting}
cd Topology/UnknotRecognition/fast
python3 -m unittest discover -s tests -v
python3 -m fastunknot recognize examples/torus_3_5.json
python3 -m fastunknot recognize examples/torus_3_5.json --no-braid
\end{lstlisting}
The paired benchmark script accepts the original baseline separately so that
paired comparisons use unchanged baseline code. The baseline module
files can be checked against the recorded Git blob hashes.

The supplied build script runs three pdfLaTeX passes to resolve the table
of contents and cross-references; the archive also contains the compiled
PDF. The release contains no claimed Lean
formalization. Its proof status is conventional mathematical proof,
supported by independent implementation checks and exact arithmetic.

\section{Conclusions}

The continuation supplies an exact complete decision backend whose cost
is linear in explicit three-braid word length at the symbol level,
and a certified polynomial extension through elementary strand descent.
It removes the need to construct Khovanov complexes on those inputs and
adds a linear writhe obstruction for general supplied braids.

For the general scanner, the scalar residue identifies multiplicities
that every exhaustive pivot strategy must retain at a fixed prefix.
The finite radical formula specifies a possible batch algorithm, while
the width--multiplicity bound names the parameter still needing control.
The geometric route likewise has explicit representation and cost
obligations.

These results improve the implementation and narrow the next research
targets. They do not establish a general quasi-polynomial complexity
bound for the delivered recognizer.

\begin{thebibliography}{99}
\small
\raggedright
\bibitem{proveit}
V. Reshetnikov and ProveIt contributors.
\emph{Topology/UnknotRecognition: implementations, proposals, and synthesis}.
Pinned commit \code{4e6fe879e5238ec0b134b6d33fdca0b8c7c1c711}, 2026.
\url{https://github.com/VladimirReshetnikov/ProveIt/tree/4e6fe879e5238ec0b134b6d33fdca0b8c7c1c711/Topology/UnknotRecognition}.

\bibitem{lackenbytalk}
M. Lackenby.
\emph{Unknot recognition in quasi-polynomial time}. February 2021 slides;
related Oxford lecture announcement, 2021.
\url{https://www.maths.ox.ac.uk/node/38304};
\url{https://people.maths.ox.ac.uk/lackenby/quasipolynomial-talk.pdf}.

\bibitem{lackenby2026}
M. Lackenby.
\emph{Incompressible surfaces, hierarchies and unknot recognition}.
arXiv:2607.23350v1, 2026.
\url{https://arxiv.org/abs/2607.23350}.

\bibitem{kelomaki2026}
T. Kelom\"aki and D. Sch\"utz.
\emph{On computational complexity of Khovanov homology}.
arXiv:2601.02119v1, 2026.
\url{https://arxiv.org/abs/2601.02119}.

\bibitem{birmanmenasco}
J. S. Birman and W. W. Menasco.
A note on closed 3-braids.
\emph{Communications in Contemporary Mathematics}
10, suppl.\ 1 (2008), 1033--1047.
\url{https://doi.org/10.1142/S0219199708003150};
\url{https://arxiv.org/abs/0802.1072}.

\bibitem{birmanhilden}
J. S. Birman and H. M. Hilden.
Heegaard splittings of branched coverings of \(S^3\).
\emph{Transactions of the American Mathematical Society}
213 (1975), 315--352.
\url{https://www.math.columbia.edu/~jb/hs_of_bcovs.pdf}.

\bibitem{datta2020}
A. Datta.
\emph{The braid group \(B_3\) in the framework of continued fractions}.
arXiv:2008.02262, 2020.
\url{https://arxiv.org/abs/2008.02262}.

\bibitem{datta2022}
A. Datta.
\emph{A novel connection between integral binary quadratic forms and knot
polynomials}. arXiv:2204.13660, 2022.
\url{https://arxiv.org/abs/2204.13660}.

\bibitem{bennequin}
D. Bennequin.
Entrelacements et \'equations de Pfaff.
\emph{Ast\'erisque} 107--108 (1983), 87--161.
\url{https://www.numdam.org/item/AST_1983__107-108__87_0/}.

\bibitem{morton}
H. R. Morton.
An irreducible 4-string braid with unknotted closure.
\emph{Mathematical Proceedings of the Cambridge Philosophical Society}
93 (1983), 259--261.
\url{https://doi.org/10.1017/S0305004100060540}.

\bibitem{birmansurvey}
J. S. Birman and T. E. Brendle.
Braids: a survey. In \emph{Handbook of Knot Theory},
Elsevier, 2005, 19--103.
\url{https://doi.org/10.1016/B978-044451452-3/50003-4};
\url{https://www.math.columbia.edu/~jb/Handbook-21.pdf}.

\bibitem{khovanovtangles}
M. Khovanov.
A functor-valued invariant of tangles.
\emph{Algebraic \& Geometric Topology} 2 (2002), 665--741.
\url{https://arxiv.org/abs/math/0103190};
\url{https://msp.org/agt/2002/2-2/agt-v2-n2-p03-p.pdf}.

\bibitem{barnatan}
D. Bar-Natan.
Fast Khovanov homology computations.
\emph{Journal of Knot Theory and Its Ramifications} 16 (2007), 243--255.
\url{https://doi.org/10.1142/S0218216507005294};
\url{https://arxiv.org/abs/math/0606318}.

\bibitem{kronheimer}
P. B. Kronheimer and T. S. Mrowka.
Khovanov homology is an unknot-detector.
\emph{Publications Math\'ematiques de l'IH\'ES} 113 (2011), 97--208.
\url{https://doi.org/10.1007/s10240-010-0030-y};
\url{https://arxiv.org/abs/1005.4346}.

\bibitem{crainic}
M. Crainic.
\emph{On the perturbation lemma, and deformations}.
arXiv:math/0403266, 2004.
\url{https://arxiv.org/abs/math/0403266}.
\end{thebibliography}
\end{document}
